\ssscp{Vowel in penultimate syllable}{14-04-03-02}
The penultimate vowel in most cognates of **mantəR is \it{a}.\footnote{
		The back vowels in the \tsc{Babar} languages
		are regular developments of *a /{\#}N{\gap} (see {\trf{04-55}}).}
Several languages in different subgroups have penultimate \it{e},
and one language, East \ili{Kola} (\tsc{Aru}), has penultimate \it{o}.
Penultimate \it{e} occurs in \ili{Aru} languages (except East \ili{Kola}),
\ili{Banda}, \ili{Fordata}, \ili{Kei}, and \ili{Teor} (\tsc{Seram-Tanimbar-Bomberai}).
These languages are all contiguous\footnote{
		Recall that \ili{Banda} has been spoken on \ili{Kei} Besar since soon after 1621
		and has not been spoken on the \ili{Banda} islands for centuries (\crf{ch:Ban}).}
and point to a local irregular development of penultimate *a
> **ə; **mantəR > **məntəR.
Penultimate *ə accounts for the vowels in all these languages.
(see \trf{14-08} for examples of penultimate *ə > \it{e} in \ili{Aru}
languages,\footnote{
		Goda-goda, with \it{modah} ‘cuscus’, is a village in the West \ili{Kola} area,
		and one other example of penultimate *ə > \it{o} is attested
		after a labial: *bətak > East \ili{Kola} -\it{ɸota} ‘split in two’.
		Otherwise the usual reflex of penultimate *ə is \it{e}.}
\srf{12-05} for penultimate *ə > \it{e} in \ili{Banda},
and reflexes of *dəŋəR [\trf{11-04}],
*qaləjaw [\trf{11-07}], and *baqəRu [\trf{11-28}] for *ə > \it{e}
in \ili{Fordata} and \ili{Kei}.) 

Waima{\Q}a \it{mide} and \ili{Tetun} \it{meda} also have irregular
penultimate vowels which cannot be reconciled with either *a or *ə.
\ili{Tetun} has penultimate *ə > \it{o}
and Waima{\Q}a has penultimate *ə > \it{e}.
(Both final vowels of these terms are regular for *ə in final
syllables.) 
To summarise, penultimate *a can account for most of the data straightforwardly.
Some forms have penultimate \it{e} for which we posit a local
irregular development of *a > **ə.
That this development is geographically restricted
and occurs in part(s) of three separate subgroups indicates that
**mantəR spread by contact to some extent.
\citet[264]{Schapper2011} suggests that \ili{Aru} reflexes of **mantəR
were introduced from \ili{Kei}.